% vim: set textwidth=78 autoindent:

\section{Externe Python Plugins verwenden}\label{sec:external_plugins}\index{plugins}

% when the revision of a section has been finalized, 
% comment out the following line:
% \updatedisclaimer

Externe QGIS Plugins sind in der Sprache Python geschrieben. Sie sind auf
einem Server in einem offiziellen, moderierten Repository abgelegt und werden
durch den jeweiligen Autor gepflegt. Tabelle \ref{tab:external_plugins} zeigt
eine Liste der aktuellen offiziellen Plugins mit einer kurzen Beschreibung.
\footnote{Auch Updates von Kern-Plugins wie dem MapServer Export k�nnen im
offiziellen Repository bereitgestellt werden.}
\footnote{Der Python Plugin Installer ist auch ein externes Plugin, aber
gleichzeitig auch Teil des QGIS Quellcodes und damit automatisch wie ein
Kern-Plugin im Plugin Manager ausw�hlbar (siehe
Kapitel~\ref{sec:load_external_plugin}).}

Als dieses Handbuch ver�ffentlich wurde, waren das offizielle, moderierte
externe Repository noch nicht komplett umgesetzt. Detailierte Informationen
zur Benutzung, zum Autor und weitere Hinweise werden jeweils mit dem Plugin
selbst bereitgestellt und sind nicht Teil dieser Dokumentation. 

Eine wahrscheinlich aktuellere Liste der moderierten externen Plugins finden
Sie im QGIS Official Repository im Men�
\dropmenuopttwo{plugin_installer}{Python Plugins herunterladen} und unter der
URL: \url{http://qgis.osgeo.org/download/plugins.html}. 

\begin{table}[H]
\centering
\caption{Moderierte externe QGIS Python Plugins}\label{tab:external_plugins}\medskip
\small
 \begin{tabular}{|l|l|p{4in}|}
\hline \textbf{Icon} & \textbf{externes Plugin} & \textbf{Beschreibung}\\
\hline
\includegraphics[width=0.7cm]{zoom2point_icon}
 & Zoom To Point \index{plugins!Zoom To Point} & Zoomed zu einer Koordinate,
die im Dialog angegeben wurde. Ein Zoom-Level kann zur Kontrolle des
Kartenausschnitts angegeben werden. \\
\hline
\end{tabular}
\end{table}

Weitere Informationen zur Installation von externen Python Plugins finden Sie
in Kapitel~\ref{sec:load_external_plugin}.

\minisec{User-Contributed Python Plugin Repository and author repositories}

Neben den moderierten externen Plugins gibt es auch ein weiteres
inoffizielles Repository. Es enth�lt eine Reihe interessanter Plugins, die
aber noch nicht stabil genug sind, um in das offizielle Repository zu kommen.
Ausserdem gibt es auch noch Repositories einzelner Anwender, die ihre Plugins
unabh�ngig vom QGIS Projekt pflegen. 

Um ein inoffizielles oder Anwender-Repository hinzuzuf�gen, �ffnen Sie den
Plugin Installer im Men� \mainmenuopt{Plugins} >
\dropmenuopttwo{plugin_installer}{Python Plugins herunterladen}. Nun wechseln
Sie in den Reiter \tab{Repositorys} und klicken auf den Knopf \button{Plugin
Repositorys hinzuf�gen}. Sie k�nnen bereits hinzugef�gte Plugin Repositorys
auch deaktivieren oder ganz aus der Liste l�schen.

\begin{Tip} \caption{\textsc{Weitere externe Plugin Repositorys hinzuf�gen}}
\qgistip{Neben dem offiziellen QGIS Plugin Repository k�nnen Sie auch
weiteres Repository hinzuf�gen. Dazu w�hlen Sie den Reiter \tab{Repositorys}
im Python Plugins Installer.
}
\end{Tip}

